% Draft subsection for the Results section of main.tex.
% Suggested placement: after the 3.1 subsection
% (\subsection{Quantitative measurements reproduce annotated facial phenotypes}).
%
% All numbers recomputed 2026-08-25 on the 120-column extractor in
% canonical-depth pose mode (mode C, docs/FEATURE_AUDIT.md 5.7). The previous
% values were computed with the per-image-z correction (arm D) and are kept in
% facekit_analysis/analysis_armD/. Sources:
%   /vast/projects/kai/multimodal-machine-learn/hongzhuo/facekit_analysis/analysis/
%   - summary.csv                              (per-disease sign match, rho, r)
%   - per_disease/*.csv                        (within-disease Cohen's d)
%   - failure_modes/sign_all.csv               (265 disease-feature pairs)
%   - failure_modes/variance_ratio_per_region.csv
%
% TODO before merging into main.tex:
%   - figures available but not yet copied into manuscript/image/:
%       analysis/summary_sign_match.png, analysis/summary_correlation.png,
%       analysis/failure_modes/sign_mismatch_by_region.png,
%       analysis/failure_modes/variance_ratio.png
%   - add a PDIDB citation/accession once the dataset record is public

\subsection{Synthetic images reproduce cohort-level geometric phenotypes}

We next asked whether the synthetic faces produced by the trained generators
carry the same geometric phenotype signal as the real photographs they were
trained to imitate. For ten of the rare-disease cohorts we compared 1,727
synthetic images against the 1,647 real images of the same syndromes, passing
both sets through the identical FaceKit pipeline and comparing them on the same
120 measurements. For each syndrome the comparison was restricted to the
measurements its HPO annotations map to, giving 265 syndrome--measurement pairs
in total.

Within each syndrome, the synthetic and real distributions of a measurement were
close. The median absolute Cohen's $d$ between the two sets ranged from 0.048
(Williams--Beuren syndrome) to 0.149 (Noonan syndrome), with an overall
median of 0.075; 84.5\% of all pairs fell below $|d| = 0.2$ and 99.6\% below
$|d| = 0.5$. A generator therefore reproduces not only the appearance of a
cohort but the location and spread of the individual measurements taken from it.

A more demanding test is whether the synthetic images preserve what
distinguishes one syndrome from the others, rather than only reproducing each
cohort in isolation. For every syndrome we computed the effect of the focal
cohort against the nine remaining cohorts pooled, separately in the synthetic
and in the real images, and compared the two sets of effects. The direction of
the effect agreed between synthetic and real for a median of 85.6\% of a
syndrome's measurements, and the effect sizes themselves tracked closely
(median Spearman $\rho = 0.912$, median Pearson $r = 0.921$). Agreement was
weakest for KBG syndrome, Cornelia de Lange syndrome and Noonan syndrome, which
tied at 80.0\% of their 35, 25 and 20 measurements, and complete for Angelman
syndrome, although the latter is mapped to only three measurements.

Across all 265 pairs, 38 (14.3\%) disagreed in sign. These disagreements were
concentrated in the eye region (23 of 126 pairs, 18.3\%) and the eyebrow region
(5 of 30, 16.7\%), and were absent from the philtrum. They were
also almost entirely confined to measurements that barely separate the cohorts
in the real images to begin with: the median $|d|$ measured in real images was
0.047 for the disagreeing pairs against 0.246 for the agreeing ones, and 92\% of
the disagreements occurred where the real effect was below $|d| = 0.2$. Where a
measurement carries no discriminative signal in the real cohort, its sign is not
determined, so these cases indicate an absence of signal to reproduce rather
than a signal reproduced incorrectly. The eyebrow region is the exception worth
naming: KBG syndrome, the weakest cohort overall, disagrees most on inter-eyebrow
distance ($d = +0.228$ synthetic against $d = -0.113$ real), the same measurement
that failed for synophrys above, and for the same reason --- the phenotype lives
in hair rather than in shape, and neither the landmark mesh nor a generator
trained through it represents it.

The one systematic difference between synthetic and real images is dispersion.
The variance of a measurement in the synthetic set was consistently smaller than
in the real set, with a median ratio of 0.87 across all pairs and regional
medians from 0.77 in the chin and jaw and 0.78 in the eyebrow region to 0.95 in
the philtrum. Synthetic cohorts are therefore narrower than the cohorts they were
trained on: the generators recover where a syndrome sits in measurement space
and how it differs from other syndromes, but under-represent the variation
within it. This should be taken into account when synthetic images are used to
augment training data, since a model trained on them will see a narrower range
of facial morphology than the real population presents.
